\chapter{Espacios de Hilbert}\label{ch:b3:hilbert}

Un \omterm{def:b3:hilbert:inner}{espacio de Hilbert} es un espacio de
Banach cuya norma procede de un \omterm{def:b3:hilbert:inner}{producto
escalar} --- y esa única estructura adicional restaura, en dimensión
infinita, casi toda la geometría euclídea: existen las proyecciones
ortogonales, toda funcional
\omterm{def:b3:topology:continuity}{continua} es un
\omterm{def:b3:hilbert:inner}{producto escalar} contra un vector fijo
(Riesz), y las bases ortonormales desarrollan cada vector en una serie
convergente con contabilidad pitagórica
(\omterm{thm:b3:hilbert:parseval}{Parseval}). El punto culminante del
capítulo es una deuda saldada: el sistema trigonométrico es una base
ortonormal de $L^2$, de modo que la identidad de
\omterm{thm:b3:hilbert:parseval}{Parseval} vale para \emph{toda} función
de cuadrado \omterm{def:b3:lebesgue:l1}{integrable} --- el enunciado que
en segundo año solo pudo demostrarse para funciones $\mathcal C^1$ a
trozos. Terminamos con Lax--Milgram, el lema caballo de batalla del
enfoque variacional de las ecuaciones diferenciales.

En todo el capítulo, $H$ es un espacio vectorial sobre $K = \R$ o $\C$.

\section{Productos escalares; el teorema de la proyección}

\begin{definition}\label{def:b3:hilbert:inner}
Un \emph{producto escalar}\index{producto escalar} es una aplicación
$\langle
\cdot,\cdot\rangle \colon H\times H \to K$, lineal en la segunda
variable, con $\langle y, x\rangle =
\overline{\langle x, y\rangle}$ y $\langle x, x\rangle > 0$ para
$x \neq 0$. Induce la norma $\norm x = \langle x,
x\rangle^{1/2}$, la \emph{desigualdad de Cauchy--Schwarz}
$\abs{\langle x, y\rangle} \leq \norm x\norm y$ (la demostración de
segundo año --- el discriminante --- no cambia) y la \emph{ley del
paralelogramo}
\[
\norm{x + y}^2 + \norm{x - y}^2 = 2\norm x^2 + 2\norm y^2 .
\]
Un \emph{espacio de Hilbert}\index{espacio de Hilbert} es un espacio con
producto escalar \omterm{def:b3:complete:complete}{completo} para esa
norma. Ejemplos: $\ell^2$ (el~\cref{pb:b3:banach:1}) y, el fundamental,
$L^2(\mu)$ con $\langle f, g\rangle = \int\bar fg\,\dd\mu$ ---
\omterm{def:b3:complete:complete}{completo} por Riesz--Fischer
(el~\cref{thm:b3:lp:complete}); el producto escalar es finito por
Cauchy--Schwarz ($= $ Hölder en
$p = q = 2$).
\end{definition}

\begin{theorem}[Proyección sobre un convexo cerrado]
\label{thm:b3:hilbert:projection}
Sean $C \neq \varnothing$ un subconjunto \emph{convexo} cerrado del
\omterm{def:b3:hilbert:inner}{espacio de Hilbert} $H$ y $x \in H$.
Existe un único $p_C(x)
\in C$ con\index{teorema de la proyección}
\[
\norm{x - p_C(x)} = d(x, C),
\]
caracterizado por: $\operatorname{Re}\langle x - p_C(x),\ c -
p_C(x)\rangle \leq 0$ para todo $c \in C$. La aplicación $p_C$ es
$1$-lipschitziana.
\end{theorem}

\begin{proof}
Sean $d = d(x, C)$ y $(c_n) \subseteq C$ con $\norm{x - c_n}
\to d$. Paralelogramo sobre $x - c_n$ y $x - c_m$:
\[
\norm{c_n - c_m}^2 = 2\norm{x - c_n}^2 + 2\norm{x - c_m}^2 -
4\,\bigl\|x - \tfrac{c_n + c_m}2\bigr\|^2
\leq 2\norm{x{-}c_n}^2 + 2\norm{x{-}c_m}^2 - 4d^2
\]
(la convexidad pone el punto medio en $C$): el miembro derecho tiende a
$0$, luego $(c_n)$ es de Cauchy, y su límite $p \in C$ (cerrado) alcanza
$d$. Unicidad: dos minimizadores dan, por la misma identidad,
$\norm{p - p'}^2 \leq 2d^2 + 2d^2 - 4d^2 = 0$.

Caracterización: para $c \in C$, $t \in \intoc01$, el vector
$p + t(c - p) \in C$, luego
\[
d^2 \leq \norm{x - p - t(c-p)}^2
= d^2 - 2t\operatorname{Re}\langle x - p, c - p\rangle +
t^2\norm{c-p}^2 ;
\]
divídase por $t \to 0^+$: $\operatorname{Re}\langle x - p, c -
p\rangle \leq 0$. Recíprocamente, esta desigualdad da $\norm{x -
c}^2 = \norm{x - p}^2 - 2\operatorname{Re}\langle x - p, c -
p\rangle + \norm{p - c}^2 \geq \norm{x-p}^2$. Lipschitz: para $x, y$ con
proyecciones $p, q$, súmense las dos desigualdades variacionales (con
$c = q$ y con $c = p$, respectivamente):
$\operatorname{Re}\langle x - y - (p - q), p - q\rangle \geq
0$, luego $\norm{p - q}^2 \leq \operatorname{Re}\langle x - y, p -
q\rangle \leq \norm{x - y}\norm{p - q}$.
\end{proof}

\begin{theorem}[Descomposición ortogonal]
\label{thm:b3:hilbert:decomposition}
Sea $F$ un \emph{subespacio cerrado} de $H$. Entonces $p_F$ es lineal,
$x - p_F(x) \perp F$ para todo $x$, y\index{complemento ortogonal}
\[
H = F \oplus F^\perp,
\qquad F^\perp = \{y : \langle y, f\rangle = 0\ \forall f\in
F\},
\qquad (F^\perp)^\perp = F .
\]
Para un subespacio general, $(F^\perp)^\perp = \bar F$; en particular,
$F$ es denso si y solo si $F^\perp = \{0\}$.
\end{theorem}

\begin{proof}
Para un subespacio, la caracterización variacional con $c =
p_F(x) \pm f$ ($f \in F$, con ambos signos, y $\iu f$ en el caso
complejo) obliga a $\langle x - p_F(x), f\rangle = 0$: el residuo es
ortogonal a $F$. Descomposición $x = p_F(x) + (x
- p_F(x))$ con $F \cap F^\perp = \{0\}$ ($\langle y, y
\rangle = 0$); la linealidad de $p_F$ se sigue de la unicidad de tales
descomposiciones (ambos miembros son lineales en ellas). $(F^\perp)
^\perp \supseteq F$ siempre; recíprocamente, si $x \perp F^\perp$,
escríbase $x = f + g$: $g = x - f \in F^\perp$ y $\langle g,
g\rangle = \langle x, g\rangle - \langle f, g\rangle = 0$: $x
= f \in F$. Para un subespacio general $F$: $F^\perp = \bar
F^{\,\perp}$ (\omterm{def:b3:topology:continuity}{continuidad} del
\omterm{def:b3:hilbert:inner}{producto escalar}), luego
$(F^\perp)^\perp = \bar F$ por el caso cerrado; y la
\omterm{ex:b3:lebesgue:gamma}{densidad} equivale a $\bar F = H$, que
equivale a $F^\perp = 0$.
\end{proof}

\begin{theorem}[Representación de Riesz]\label{thm:b3:hilbert:riesz}
Para toda funcional lineal \omterm{def:b3:topology:continuity}{continua}
$\varphi \in H'$ existe un único $a \in H$ con\index{teorema de
representación de Riesz}
\[
\varphi(x) = \langle a, x\rangle \quad (x \in H),
\qquad \norm\varphi_{H'} = \norm a .
\]
\end{theorem}

\begin{proof}
Si $\varphi = 0$: $a = 0$. En caso contrario, $F = \ker\varphi$ es un
subespacio propio cerrado; tómese $u \in F^\perp$, $\norm u = 1$
(el~\cref{thm:b3:hilbert:decomposition}: $F^\perp \neq 0$ porque
$F \neq H$). Para todo $x$, el vector $\varphi(x)u -
\varphi(u)x \in \ker\varphi$ y, por tanto, $\perp u$:
\[
0 = \langle u, \varphi(x)u - \varphi(u)x\rangle
= \varphi(x) - \varphi(u)\langle u, x\rangle :
\qquad \varphi(x) = \langle
\overline{\varphi(u)}\,u,\ x\rangle .
\]
Así que $a = \overline{\varphi(u)}u$ sirve. Unicidad: $\langle a
- a', x\rangle = 0$ para todo $x$; evalúese en $x = a - a'$. Normas:
$\abs{\varphi(x)} \leq \norm a\norm x$ (Cauchy--Schwarz), con igualdad
en $x = a$.
\end{proof}

\begin{example}[Una proyección, calculada hasta el final]
\label{ex:b3:hilbert:projectionexample}
En $H = L^2(\intcc01)$, ¿cuál es la mejor aproximación de $f(x) = x^2$
por una función afín? El subespacio $F =
\operatorname{Vect}(1, x)$ es cerrado (de dimensión finita), y
$p_F(f) = a + bx$ queda caracterizada por la ortogonalidad del residuo a
$1$ y a $x$:
\[
\int_0^1(x^2 - a - bx)\,\dd x = 0,
\qquad
\int_0^1x\,(x^2 - a - bx)\,\dd x = 0,
\]
es decir, $\frac13 = a + \frac b2$ y $\frac14 = \frac a2 +
\frac b3$: $a = -\frac16$, $b = 1$. Luego $p_F(x^2) = x -
\frac16$, y el error es
\[
d(f, F)^2 = \int_0^1\Bigl(x^2 - x + \frac16\Bigr)^2\dd x =
\frac1{180},
\qquad d(f, F) = \frac1{6\sqrt5} .
\]
Dos observaciones que conviene interiorizar. Primera: el cálculo no es
más que un sistema lineal $2\times2$ --- las \emph{ecuaciones normales};
para la base de los monomios, su matriz $\bigl(\frac1{i+j+1}\bigr)$ es
la célebremente mal condicionada matriz de Hilbert, y ortogonalizar
primero (polinomios de Legendre, el~\cref{pb:b3:hilbert:1}) es el
remedio. Segunda: la mejor aproximación \emph{uniforme} de $x^2$ por
funciones afines es otra ($x - \frac18$, por equioscilación): cada norma
tiene su propia geometría, y solo la hilbertiana responde con un sistema
lineal.
\end{example}

\section{Bases ortonormales}

\begin{definition}\label{def:b3:hilbert:onb}
Una familia $(e_i)_{i\in I}$ es \emph{ortonormal} si $\langle
e_i, e_j\rangle = \delta_{ij}$, y una \emph{base hilbertiana}\index{base
hilbertiana} (base ortonormal) si, además, sus combinaciones lineales
finitas son densas en $H$ (la familia es \emph{total}). Trataremos el
caso numerable $I = \N$ que, por Gram--Schmidt, cubre todo $H$
\emph{separable}
(\cref{prop:b3:hilbert:gramschmidt}).
\end{definition}

\begin{theorem}[Bessel, Parseval]\label{thm:b3:hilbert:parseval}
Sea $(e_n)_{n\in\N}$ ortonormal en $H$, y $c_n(x) =
\langle e_n, x\rangle$.
\begin{enumerate}
  \item (Bessel) $\sum_n\abs{c_n(x)}^2 \leq \norm x^2$, y la serie
  $\sum_nc_n(x)e_n$ converge en $H$, con suma $p_F(x)$,
  $F = \overline{\operatorname{Vect}}(e_n)$. \item Son equivalentes: (i)
  $(e_n)$ es una \omterm{def:b3:hilbert:onb}{base hilbertiana}; (ii)
  $x = \sum_nc_n(x)e_n$ para todo $x$; (iii)
  \emph{Parseval}\index{identidad de Parseval}: $\norm x^2
        = \sum_n\abs{c_n(x)}^2$ para todo $x$; (iv) el único vector
  ortogonal a todos los $e_n$ es $0$. \item Si $(e_n)$ es una
  \omterm{def:b3:hilbert:onb}{base hilbertiana}, $x \mapsto (c_n(x))_n$
  es un isomorfismo isométrico $H \to \ell^2$ (\emph{todo}
  \omterm{def:b3:hilbert:inner}{espacio de Hilbert} separable de
  dimensión infinita «es» $\ell^2$), y $\langle x, y\rangle =
        \sum_n\overline{c_n(x)}c_n(y)$.
\end{enumerate}
\end{theorem}

\begin{proof}
(1) Para $N$ finito: $x - \sum_{n\leq N}c_ne_n \perp e_k$ ($k
\leq N$), de modo que Pitágoras da $\norm x^2 = \sum_{n\leq
N}\abs{c_n}^2 + \norm{x - \sum_{n\leq N}c_ne_n}^2$: Bessel. Las sumas
parciales $S_N = \sum_{n\leq N}c_ne_n$ son de Cauchy:
$\norm{S_N - S_M}^2 = \sum_{M<n\leq N}\abs{c_n}^2$, cola de una serie
convergente; el límite está en $F$, y $x - \lim S_N
\perp$ cada $e_k$ (\omterm{def:b3:topology:continuity}{continuidad}),
luego $\perp F$: por la unicidad de la descomposición ortogonal,
$\lim S_N = p_F(x)$.

(2) (i)$\Rightarrow$(ii): $F = H$, luego $p_F = \mathrm{id}$.
(ii)$\Rightarrow$(iii): Pitágoras en el límite ($\norm{S_N}^2
= \sum_{n \leq N}\abs{c_n}^2 \to \norm x^2$). (iii)$\Rightarrow$(iv):
$x \perp$ para todos los $e_n$ da $\norm x^2 =
0$. (iv)$\Rightarrow$(i): $F^\perp = \{0\}$ (ser ortogonal a todos los
$e_n$ es serlo a $F$), de modo que $F$ es denso por
el~\cref{thm:b3:hilbert:decomposition}; pero $F$, una
\omterm{def:b3:topology:interior}{clausura}, ya es cerrado: $F = H$.

(3) La aplicación es lineal, isométrica por (iii) (luego inyectiva) y
sobreyectiva: dado $(c_n) \in \ell^2$, la serie $\sum
c_ne_n$ converge (de Cauchy como en (1)) a una preimagen. La fórmula del
\omterm{def:b3:hilbert:inner}{producto escalar} es la polarización a
partir de (iii), o un cálculo directo de límites.
\end{proof}

\begin{proposition}[Gram--Schmidt]\label{prop:b3:hilbert:gramschmidt}
Sea $(x_n)$ una sucesión linealmente independiente. Poniendo
inductivamente $\tilde e_n = x_n - \sum_{k<n}\langle e_k,
x_n\rangle e_k$ y $e_n = \tilde e_n/\norm{\tilde e_n}$ se obtiene una
familia ortonormal $(e_n)$ con las mismas envolturas finitas:
$\operatorname{Vect}(e_1, \dots, e_n) = \operatorname{Vect}
(x_1, \dots, x_n)$. En consecuencia, todo
\omterm{def:b3:hilbert:inner}{espacio de Hilbert} separable (con un
subconjunto denso numerable) tiene una \omterm{def:b3:hilbert:onb}{base
hilbertiana}.\index{ortonormalización de Gram--Schmidt}
\end{proposition}

\begin{proof}
Inducción: $\tilde e_n \perp e_k$ ($k < n$) por construcción, y
$\tilde e_n \neq 0$ por independencia; las envolturas coinciden en cada
etapa (cambio de base triangular). Para $H$ separable: de una sucesión
densa extráigase una subfamilia linealmente independiente con envoltura
densa (descártese cada vector que esté en la envoltura de los anteriores
--- la envoltura no cambia) y ortonormalícese: el resultado es total.
\end{proof}

\begin{theorem}[El sistema trigonométrico; Parseval, por fin]
\label{thm:b3:hilbert:fourier}
En $L^2(\intcc{-\pi}\pi)$ con $\langle f, g\rangle =
\frac1{2\pi}\int_{-\pi}^\pi \bar fg$, la familia $e_n(t) =
\eu^{\iu nt}$, $n \in \Z$, es una \omterm{def:b3:hilbert:onb}{base
hilbertiana}. En consecuencia, para \emph{toda} $f \in L^2$ --- en
particular, para toda $f$ $2\pi$-periódica
\omterm{def:b3:topology:continuity}{continua} a trozos --- con $c_n(f) =
\frac1{2\pi}\int_{-\pi}^{\pi}f(t)\eu^{-\iu nt}\dd t$:
\[
f = \sum_{n\in\Z}c_n(f)\,\eu^{\iu nt} \ \ \text{en } L^2,
\qquad
\frac1{2\pi}\int_{-\pi}^{\pi}\abs f^2 =
\sum_{n\in\Z}\abs{c_n(f)}^2 .
\]
Esto demuestra, con toda generalidad, la identidad de
\omterm{thm:b3:hilbert:parseval}{Parseval} que segundo año
admitió.\index{identidad de Parseval (general)}
\end{theorem}

\begin{proof}
La ortonormalidad es un cálculo directo (segundo año). Totalidad: sea
$f \in L^2$ $\perp$ a todos los $e_n$, es decir, con todos los
coeficientes de Fourier nulos. Las funciones $2\pi$-periódicas
\omterm{def:b3:topology:continuity}{continuas} son densas en
$L^2(\intcc{-\pi}\pi)$: en efecto, $\mathcal
C_c(\intoo{-\pi}\pi)$ es denso (el~\cref{thm:b3:lp:density}(2)) y tales
funciones se extienden periódica y
\omterm{def:b3:topology:continuity}{continuamente}. Los polinomios
trigonométricos son $\norm\cdot_\infty$-densos entre las funciones
periódicas \omterm{def:b3:topology:continuity}{continuas}
(Stone--Weierstrass, el~\cref{cor:b3:complete:weierstrass}(c)), y
$\norm\cdot_2 \leq \norm\cdot_\infty$: los polinomios trigonométricos
son densos en $L^2$. Pero $f \perp$ a todo polinomio trigonométrico y,
por tanto, $f \perp$ a un subespacio denso: $f
\in (\text{denso})^\perp = \{0\}$
(el~\cref{thm:b3:hilbert:decomposition}). El criterio (iv)
del~\cref{thm:b3:hilbert:parseval} concluye; y (ii) y (iii) se traducen
en la fórmula mostrada (reindexando el conjunto numerable $\Z$; la serie
doblemente infinita converge incondicionalmente --- las sumas parciales
sobre cualquier familia exhaustiva convergen, por el argumento de la
cola $\ell^2$).
\end{proof}

\begin{theorem}[Lax--Milgram]\label{thm:b3:hilbert:laxmilgram}
Sean $H$ un \omterm{def:b3:hilbert:inner}{espacio de Hilbert} real y
$a \colon H\times H \to
\R$ bilineal, \emph{\omterm{def:b3:topology:continuity}{continua}}
($\abs{a(u,v)} \leq M\norm
u\norm v$) y \emph{coerciva} ($a(u, u) \geq \alpha\norm u^2$,
$\alpha > 0$). Entonces, para toda $\varphi \in H'$ existe un único
$u \in H$ con\index{teorema de Lax--Milgram}
\[
a(u, v) = \varphi(v) \qquad \text{para todos } v \in H .
\]
\end{theorem}

\begin{proof}
Para $u$ fijo, $v \mapsto a(u, v)$ es lineal
\omterm{def:b3:topology:continuity}{continua}: Riesz da un único
$Au \in H$ con $a(u,v) = \langle Au,
v\rangle$; $A$ es lineal con $\norm{Au} \leq M\norm u$ (unicidad de los
representantes, y después la cota). Coercividad:
$\alpha\norm u^2 \leq a(u,u) = \langle Au, u\rangle \leq
\norm{Au}\norm u$, luego $\norm{Au} \geq \alpha\norm u$: $A$ es
inyectiva con imagen cerrada (una sucesión de Cauchy de imágenes $Au_n$
obliga a que $u_n$ sea de Cauchy). La imagen es densa: $w \perp
\operatorname{im}A$ da $0 = \langle Aw, w\rangle \geq
\alpha\norm w^2$. Cerrada y densa: $A$ es biyectiva. Dada $\varphi$, sea
$f$ su representante (Riesz) y $u = A^{-1}f$:
$a(u, v) = \langle f, v\rangle = \varphi(v)$, de manera única
($a(u - u', \cdot) = 0$ y coercividad).
\end{proof}

\begin{remark}\label{rem:b3:hilbert:variational}
Cuando $a$ es simétrica, la solución de Lax--Milgram es el único
minimizador de la \emph{energía} $J(v) = \frac12a(v,v) -
\varphi(v)$ (el~\cref{exo:b3:hilbert:9}): la existencia de soluciones de
problemas variacionales de un solo golpe. Aplicado a espacios de
funciones adecuados (los espacios de Sóbolev de un curso posterior),
esto resuelve problemas de contorno para ecuaciones diferenciales --- la
puerta de entrada moderna a las ecuaciones en derivadas parciales.
\end{remark}

\section{Ejercicios}

\begin{exercise}[$\star$]\label{exo:b3:hilbert:1}
(a) Demostrar las identidades de polarización (caso real: $4\langle x,
y\rangle = \norm{x+y}^2 - \norm{x-y}^2$; caso complejo: la versión de
cuatro términos). (b) Demostrar que $\norm\cdot_1$ sobre $L^1(\intcc01)$
y $\norm\cdot_\infty$ sobre $\mathcal C(\intcc01)$ violan la ley del
paralelogramo: estas normas no proceden de ningún
\omterm{def:b3:hilbert:inner}{producto escalar}.
\end{exercise}

\begin{exercise}[$\star$]\label{exo:b3:hilbert:2}
En $H = L^2(\intcc01)$ (real): (a) calcúlese la proyección de $f$ sobre
el subespacio de las funciones constantes e interprétese; (b) calcúlese
la proyección sobre $\{g : g = 0 \text{ en casi todo punto de }
\intcc0{1/2}\}$; (c) calcúlese
$d\bigl(x \mapsto x,\ \operatorname{Vect}(\mathbf
1)\bigr)$.
\end{exercise}

\begin{exercise}[$\star\star$]\label{exo:b3:hilbert:3}
(a) Demostrar que, para un subespacio $F$: $F$ denso $\iff$ $F^\perp =
\{0\}$, y dese un ejemplo en $\ell^2$ de un subespacio denso
\emph{propio} (de modo que $F^\perp = 0$ sin que $F = H$: el teorema de
descomposición necesita genuinamente que $F$ sea cerrado). (b) Demostrar
que si $x_n \to x$ y $y_n \to y$ en norma, entonces
$\langle x_n, y_n\rangle \to \langle x, y\rangle$, y localícense dos
lugares del capítulo donde se usó esta
\omterm{def:b3:topology:continuity}{continuidad}.
\end{exercise}

\begin{exercise}[$\star\star$]\label{exo:b3:hilbert:4}
Aplíquese Gram--Schmidt a $1, x, x^2$ en $L^2(\intcc{-1}1)$
(\omterm{def:b3:measure:lebesgueouter}{medida de Lebesgue}): obténganse
los tres primeros \emph{polinomios de Legendre} normalizados y
compruébese que coinciden con $\sqrt{n + \frac12}\,P_n$ para los
polinomios de Rodrigues $P_n$ del~\cref{pb:b3:hilbert:1}.
\end{exercise}

\begin{exercise}[$\star\star$]\label{exo:b3:hilbert:5}
Aplíquese \omterm{thm:b3:hilbert:parseval}{Parseval}
(el~\cref{thm:b3:hilbert:fourier}) a $f(t) = t$ y a $f(t) = t^2$ en
$\intcc{-\pi}\pi$ --- ahora legítimamente para estas funciones
(\omterm{def:b3:topology:continuity}{continuas}, aunque antes la
identidad exigía cuidados de tipo $\mathcal C^1$ a trozos en la
discontinuidad de empalme): recupérense
\[
\sum_{n\geq1}\frac1{n^2} = \frac{\pi^2}6,
\qquad
\sum_{n\geq1}\frac1{n^4} = \frac{\pi^4}{90} .
\]
\end{exercise}

\begin{exercise}[$\star\star$]\label{exo:b3:hilbert:6}
(a) Hallar $a \in L^2(\intcc01)$ con $\int_0^{1/2}f =
\langle a, f\rangle$ para toda $f$; calcúlese $\norm\varphi$ para esa
funcional. (b) Demostrar que la evaluación $f \mapsto f(\frac12)$,
definida sobre el subespacio $\mathcal C(\intcc01) \subseteq
L^2(\intcc01)$, \emph{no} es
\omterm{def:b3:topology:continuity}{continua} para $\norm\cdot_2$: no
existe representante de Riesz (la evaluación no es una noción de $L^2$).
\end{exercise}

\begin{exercise}[$\star\star\star$]\label{exo:b3:hilbert:7}
Sean $H$ separable con \omterm{def:b3:hilbert:onb}{base hilbertiana}
$(e_n)$, y $(x_k)$ una sucesión acotada. (a) Demostrar que alguna
subsucesión converge \emph{débilmente}: existe $x$ con
$\langle y, x_{k_j}\rangle \to \langle y,
x\rangle$ para todo $y \in H$. \emph{(Extracción diagonal sobre los
coeficientes $\langle e_n, x_k\rangle$; constrúyase $x$ mediante Bessel
y la acotación uniforme de las normas.)} (b) Demostrar que
$e_n \rightharpoonup 0$ pero $\norm{e_n} = 1$: los límites débiles
pueden perder norma. Demuéstrese $\norm x \leq
\liminf\norm{x_{k_j}}$ en (a).
\end{exercise}

\begin{exercise}[$\star\star$]\label{exo:b3:hilbert:8}
(Adjuntos) Para $T \in \mathcal L(H)$, demostrar que existe un único
$T^* \in \mathcal L(H)$ con $\langle Tx, y\rangle = \langle
x, T^*y\rangle$ (Riesz), y que $\vertiii{T^*} = \vertiii T$. Calcúlese
el adjunto del desplazamiento $S$ en $\ell^2$ y demuéstrese
$\ker T^* = (\operatorname{im}T)^\perp$ --- dedúzcase
$\overline{\operatorname{im}T} = (\ker T^*)^\perp$.
\end{exercise}

\begin{exercise}[$\star\star$]\label{exo:b3:hilbert:9}
Sea $a$ como en Lax--Milgram y, además, \emph{simétrica}. Demostrar que
$u$ resuelve $a(u, \cdot) = \varphi$ si y solo si $u$ minimiza
$J(v) = \frac12a(v, v) - \varphi(v)$, y que el mínimo se alcanza en
exactamente un punto.
\emph{(Complétese el cuadrado:
$J(u + w) - J(u) = \frac12a(w,w) \geq
\frac\alpha2\norm w^2$.)} Aplicación: vuélvase a deducir el teorema de
la proyección sobre subespacios cerrados a partir de Lax--Milgram.
\end{exercise}

\begin{exercise}[$\star\star\star$]\label{exo:b3:hilbert:10}
(El sistema de Haar) En $\intcc01$, sea $h_{0} = \mathbf 1$ y, para
$n = 2^j + k$ ($j \geq 0$, $0 \leq k < 2^j$):
\[
h_n = 2^{j/2}\Bigl(\mathbf 1_{[k2^{-j},\,(k +
\frac12)2^{-j})} - \mathbf 1_{[(k+\frac12)2^{-j},\,(k+1)2^{-j})}
\Bigr).
\]
Demostrar que $(h_n)_{n\geq0}$ es ortonormal en $L^2(\intcc01)$ y total.
\emph{(Ortogonalidad: soportes disjuntos o encajados; totalidad: las
envolturas finitas contienen todas las funciones escalonadas diádicas,
que son densas --- mediante el~\cref{thm:b3:lp:density}(1) y la
aproximación diádica de los intervalos.)} El sistema de Haar es el
antepasado de las ondículas.
\end{exercise}

\begin{exercise}[$\star\star$]\label{exo:b3:hilbert:11}
(Proyecciones ortogonales, caracterizadas) Sean $H$ un
\omterm{def:b3:hilbert:inner}{espacio de Hilbert} y
$P \in \mathcal L(H)$ con $P^2 = P$, $P \neq 0$. Demuéstrese la
equivalencia de: (i) $P$ es la proyección ortogonal sobre
$\operatorname{im}P$;
(ii) $P = P^*$ (el~\cref{exo:b3:hilbert:8}); (iii) $\vertiii P = 1$.
\emph{(Para (iii) $\Rightarrow$ (i): si algún $x \in
(\ker P)^\perp$ cumpliera $Px \neq x$, considérese $x + t(Px - x)$ --- o
directamente: para $u \in \operatorname{im}P$ y $v \in
\ker P$, desarróllese $\norm{P(u + tv)}^2 \leq \norm{u + tv}^2$ para
todo $t \in \R$ y conclúyase $\langle u, v\rangle = 0$.)} Exhíbase una
proyección no ortogonal en $\R^2$ y calcúlese su norma.
\end{exercise}

\begin{exercise}[$\star\star\star$]\label{exo:b3:hilbert:12}
(Teorema ergódico de von Neumann) Sean $U \in \mathcal L(H)$
\emph{unitario} ($U^*U = UU^* = I$), $F = \ker(U - I)$ el espacio de
puntos fijos, $P$ la proyección ortogonal sobre $F$ y
$A_n = \frac1n\sum_{k=0}^{n-1}U^k$.
(a) Demostrar $\ker(U - I) = \ker(U^* - I)$ \emph{(a partir de
$\norm{Ux - x}^2 = 2\norm x^2 - 2\operatorname{Re}\langle
Ux, x\rangle$ y de la unitariedad)} y dedúzcase
$\overline{\operatorname{im}(U - I)} = F^\perp$.
(b) Demostrar que $A_nx \to x$ para $x \in F$, y $A_nx \to 0$ para
$x \in \operatorname{im}(U - I)$ \emph{(telescopado)}, y después para
$x \in \overline{\operatorname{im}(U - I)}$ (cota uniforme
$\vertiii{A_n} \leq 1$). (c) Concluir: $A_nx \to Px$ para \emph{todo}
$x \in H$ --- las medias temporales convergen a la proyección sobre los
invariantes. (d) Detállese para $H = L^2(\R/\Z)$ y $Uf = f(\cdot +
\alpha)$ con $\alpha$ irracional: identifíquese $F$ (úsense series de
Fourier, el~\cref{thm:b3:hilbert:fourier}) y dedúzcase que
$\frac1n\sum_{k<n}f(x + k\alpha) \to \int_0^1f$ en $L^2$: la
equidistribución $L^2$ de las rotaciones irracionales.
\end{exercise}

\section{Problema: polinomios ortogonales}

\begin{problem}[{Problema de fin de semana --- Legendre, Hermite y
cuadratura de Gauss}]\label{pb:b3:hilbert:1} Sean $I \subseteq \R$ un
intervalo y $w > 0$ un \emph{peso}
\omterm{def:b3:topology:continuity}{continuo} sobre el
\omterm{def:b3:topology:interior}{interior} de $I$ tal que $\int_I
\abs t^nw(t)\dd t < \infty$ para todo $n$; trabajaremos en $H = L^2(I,
w\,\dd\lambda)$ con $\langle f, g\rangle = \int_I \bar
fg\,w$. Gram--Schmidt aplicado a $1, t, t^2, \dots$ produce los
\emph{\omterm{pb:b3:hilbert:1}{polinomios ortogonales}} $(p_n)$ para $w$
(normalización mónica: $p_n = t^n + \cdots$).\index{polinomios
ortogonales}

\medskip
\noindent\textbf{Parte I --- Teoría general.}
\begin{enumerate}
  \item Demostrar que $p_n$ es ortogonal a todo polinomio de grado $< n$
  y que $(p_0, \dots, p_n)$ es una base de $\R_n[t]$. \item (Recurrencia
  a tres términos) Demostrar que existen reales $a_n,
        b_n$ con
        \[
        p_{n+1}(t) = (t - a_n)\,p_n(t) - b_n\,p_{n-1}(t),
        \qquad b_n = \frac{\norm{p_n}^2}{\norm{p_{n-1}}^2} >
        0 .
        \]
        \emph{(Desarróllese $t\,p_n$ en la base $(p_k)_{k \leq
        n+1}$ y anúlense coeficientes por ortogonalidad, usando
        $\langle tp_n, p_k\rangle = \langle p_n,
        tp_k\rangle$.)}
  \item (Raíces) Demostrar que $p_n$ tiene $n$ raíces \emph{distintas},
  todas \omterm{def:b3:topology:interior}{interiores} a $I$. \emph{(Sean
  $t_1 < \dots < t_m$ los cambios de signo
  \omterm{def:b3:topology:interior}{interiores} de $p_n$; si $m < n$,
  evalúese $p_n$ contra $\prod_{i\leq m}(t - t_i)$ y contradígase la
  ortogonalidad.)}
\end{enumerate}

\medskip
\noindent\textbf{Parte II --- Legendre ($I = \intcc{-1}1$, $w =
1$).} Defínase $P_n(t) = \frac{1}{2^nn!}\,\frac{\dd^n}{\dd
t^n}\bigl[(t^2 - 1)^n\bigr]$ (Rodrigues).
\begin{enumerate}[resume]
  \item Demostrar que $\deg P_n = n$ con coeficiente director
  $\frac{(2n)!}{2^n(n!)^2}$ y, integrando por partes $n$ veces, que
  $\langle P_n, Q\rangle = 0$ para todo polinomio $Q$ de grado $< n$:
  los $P_n$ son (salvo normalización) los
  \omterm{pb:b3:hilbert:1}{polinomios ortogonales} para $w = 1$. \item
  Calcular $\norm{P_n}_2^2 = \frac{2}{2n+1}$ \emph{(intégrese por partes
  $n$ veces contra sí mismo y redúzcase a una integral de tipo Beta o de
  Wallis,
        \cref{exo:b3:product:8})}.
  \item Demostrar que los polinomios de Legendre normalizados forman una
  \omterm{def:b3:hilbert:onb}{base hilbertiana} de $L^2(\intcc{-1}1)$
  \emph{(Weierstrass, el~\cref{cor:b3:complete:weierstrass}, más la
  \omterm{ex:b3:lebesgue:gamma}{densidad} de $\mathcal C$ en $L^2$)}, y
  desarróllese $f(t) = \abs t$ hasta grado $2$: calcúlese la mejor
  aproximación cuadrática en $L^2$ de $\abs t$.
\end{enumerate}

\medskip
\noindent\textbf{Parte III --- Hermite ($I = \R$, $w(t) =
\eu^{-t^2}$).} Defínase $H_n(t) =
(-1)^n\eu^{t^2}\frac{\dd^n}{\dd t^n}\eu^{-t^2}$.
\begin{enumerate}[resume]
  \item Demostrar que $H_n$ es un polinomio de grado $n$ con coeficiente
  director $2^n$, que $H_{n+1} = 2tH_n -
        H_n'$ y que $\langle H_m, H_n\rangle_w =
        \delta_{mn}\,2^nn!\sqrt\pi$ \emph{(de nuevo, por partes)}. \item
  Demostrar que la familia de Hermite es total en $L^2(\R,
        \eu^{-t^2}\dd t)$, admitiendo un resultado
  del~\cref{ch:b3:fouriertransform}: si $g \in L^1(\R)$ cumple
  $\int g(t)\eu^{-\iu\xi t}\dd t = 0$ para todo $\xi$, entonces $g = 0$
  en casi todo punto. \emph{(Para $f \perp$ todos los $H_n$, es decir,
  $\perp$ todos los polinomios: demuéstrese que $z \mapsto \int
        f(t)\eu^{-t^2}\eu^{-\iu zt}\dd t$ está bien definida,
  desarróllese la exponencial en serie, justifíquese el intercambio por
  dominación y concluir que la transformada de Fourier de $f\eu^{-t^2}$
  se anula.)}
\end{enumerate}

\medskip
\noindent\textbf{Parte IV --- Cuadratura de Gauss.} Fíjese $n$, sean
$t_1 < \dots < t_n$ las raíces de $p_n$ (Parte I) y defínanse los pesos
$w_i = \int_I \ell_i(t)\,w(t)\dd t$, donde los $\ell_i$ son los
polinomios de la base de interpolación de Lagrange en los $t_i$.
\begin{enumerate}[resume]
  \item Demostrar que la fórmula de cuadratura $Q(f) = \sum_iw_if(t_i)$
  es exacta sobre todos los polinomios de grado $\leq n - 1$
  (interpolación) y, de hecho --- el milagro ---, sobre todos los de
  grado $\leq 2n - 1$: escríbase $P =
        qp_n + r$ y úsese la ortogonalidad sobre el cociente $q$. \item
  Demostrar que los pesos son positivos \emph{(aplíquese la regla a
  $\ell_i^2$, de grado $2n -
        2$)} y dedúzcase del teorema de Pólya
  (el~\cref{exo:b3:banach:9}) que la cuadratura de Gauss converge:
  $Q_n(f) \to \int_I fw$ para toda $f$
  \omterm{def:b3:topology:continuity}{continua} sobre un $I$
  \omterm{def:b3:topology:compact}{compacto}. \item Para $n = 2$,
  $I = \intcc{-1}1$, $w = 1$: calcúlense los nodos $\pm\frac1{\sqrt3}$ y
  los pesos $1, 1$, y compruébese a mano la exactitud sobre
  $1, t, t^2, t^3$. Compárese con la regla del trapecio sobre esos
  mismos dos puntos de evaluación.
\end{enumerate}

\medskip
\noindent\textbf{Parte V --- Chebyshev: los polinomios que mejor
oscilan.} Ahora $I = \intcc{-1}1$ y $w(t) =
\frac1{\sqrt{1 - t^2}}$.
\begin{enumerate}[resume]
  \item Demostrar que $T_n(\cos\theta) = \cos n\theta$ define un
  polinomio $T_n$ de grado $n$ (establézcase $T_{n+1} =
        2t\,T_n - T_{n-1}$ a partir de una identidad trigonométrica),
  con coeficiente director $2^{n-1}$ para $n \geq 1$; y que la
  sustitución $t = \cos\theta$ da
        \[
        \langle T_m, T_n\rangle_w =
        \int_0^\pi\cos m\theta\,\cos n\theta\,\dd\theta
        = 0 \ (m \neq n), \qquad
        \norm{T_0}_w^2 = \pi,\ \ \norm{T_n}_w^2 = \frac\pi2 :
        \]
        los $T_n$ son los \omterm{pb:b3:hilbert:1}{polinomios
        ortogonales} para este peso, y los desarrollos de Chebyshev
        \emph{son} series de Fourier en cosenos disfrazadas. \item
        Localícense explícitamente las $n$ raíces $t_k =
        \cos\frac{(2k-1)\pi}{2n}$ y los $n + 1$ extremos
        $s_j = \cos\frac{j\pi}n$ de $T_n$ en $\intcc{-1}1$, donde
        $T_n(s_j) = (-1)^j$: la gráfica \emph{equioscila} entre $\pm1$.
        \item (Minimax) Demostrar que, entre todos los polinomios
        \emph{mónicos} de grado $n$, el polinomio $2^{1-n}T_n$ tiene la
        menor norma del supremo en $\intcc{-1}1$, a saber $2^{1-n}$ ---
        y que es el único minimizador. \emph{(Si un $P$ mónico cumpliera
        $\sup\abs P < 2^{1-n}$, la diferencia $2^{1-n}T_n -
        P$, de grado $\leq n-1$, alternaría de signo en los $n+1$ puntos
        de equioscilación.)} \item Aplicación a la interpolación: para
        nodos $t_1 < \dots
        < t_n$ en $\intcc{-1}1$, el error de la interpolación de
        Lagrange de una función $\mathcal C^n$ involucra
        $\omega(t) = \prod_i(t - t_i)$. Demostrar que tomar las raíces
        de Chebyshev como nodos minimiza $\sup_{\intcc{-1}1}\abs\omega$,
        y dese la cota resultante $\norm{f -
        L_nf}_\infty \leq \frac{\norm{f^{(n)}}_\infty}
        {2^{n-1}\,n!}$ --- compárese con nodos equiespaciados (enúnciese
        el fenómeno de Runge como advertencia). \item Comprobar
        $\abs{T_n'(\pm1)} = n^2$ \emph{(derívese
        $T_n(\cos\theta) = \cos n\theta$ y tómense límites
        $\theta \to 0, \pi$)}: los polinomios acotados por $1$ en
        $\intcc{-1}1$ pueden tener derivada tan grande como $n^2$ en el
        borde (la desigualdad de Markov dice que no más --- solo el
        enunciado). ¿En qué parte del intervalo la cota de la derivada
        es solo $O(n)$? \item (Cuadratura de Chebyshev--Gauss) Demostrar
        que la regla de Gauss para el peso $w$ en las $n$ raíces de
        Chebyshev tiene pesos \emph{iguales} $w_i = \frac\pi n$
        \emph{(exactitud sobre $T_0, \dots, T_{n-1}$ más las sumas
        trigonométricas
        $\sum_{k=1}^n\cos\bigl(j\tfrac{(2k-1)\pi}{2n}\bigr) =
        0$ para $1 \leq j \leq n - 1$)}: la más uniforme de todas las
        cuadraturas. Escríbase para $n = 3$.
\end{enumerate}

\medskip
\noindent\textbf{Parte VI --- Christoffel--Darboux, entrelazado y la
matriz de Jacobi.} Volvamos a un peso general; $h_k = \norm{p_k}^2$
($p_k$ mónicos), $b_k = h_k/h_{k-1}$.
\begin{enumerate}[resume]
  \item (Norma mínima) Demostrar que, entre todos los polinomios
  \emph{mónicos} de grado $n$, el ortogonal $p_n$ es el único de norma
  $L^2(w)$ mínima --- identifíquese la minimización como una proyección
  ortogonal sobre $\R_{n-1}[t]$ (el~\cref{thm:b3:hilbert:projection}, o
  la proyección de dimensión finita de segundo año). La propiedad
  minimax de la pregunta 14 es el mismo enunciado con $L^\infty$ en
  lugar de $L^2$: el mismo héroe, dos normas. \item
  (Christoffel--Darboux) Demostrar, por inducción sobre $n$ usando la
  recurrencia a tres términos, la identidad
        \[
        \sum_{k=0}^{n}\frac{p_k(x)\,p_k(y)}{h_k}
        = \frac{p_{n+1}(x)\,p_n(y) -
        p_n(x)\,p_{n+1}(y)}{h_n\,(x - y)}
        \qquad (x \neq y),
        \]
        y su forma confluente ($y \to x$):
        $\sum_{k\leq n}\frac{p_k(x)^2}{h_k} =
        \frac{p_{n+1}'(x)p_n(x) - p_n'(x)p_{n+1}(x)}{h_n}$.
  \item Deducir que $p_n$ y $p_{n+1}$ no tienen raíces comunes, y que en
  toda raíz $x_0$ de $p_{n+1}$: $p_n(x_0)\,p_{n+1}'(x_0) > 0$. Concluir
  el \emph{entrelazado} de las raíces: entre dos raíces consecutivas de
  $p_{n+1}$ hay exactamente una raíz de $p_n$. \item (Matriz de Jacobi)
  Sea $J_n$ la matriz $n\times n$ tridiagonal simétrica con diagonal
  $a_0,
        \dots, a_{n-1}$ y entradas fuera de la diagonal $\sqrt{b_1},
        \dots, \sqrt{b_{n-1}}$. Demostrar por inducción que
  $\det(tI_n - J_n) = p_n(t)$, de modo que las raíces de $p_n$ son los
  valores propios de una matriz real simétrica --- lo que vuelve a
  demostrar en una línea que son reales y (con el entrelazado anterior)
  ata los \omterm{pb:b3:hilbert:1}{polinomios ortogonales} al mundo
  espectral del~\cref{ch:b3:spectral}. \item (Síntesis) Móntese el
  diccionario de las tres familias clásicas (Legendre, Hermite,
  Chebyshev): intervalo, peso, fórmula definitoria, recurrencia a tres
  términos, norma y hábitat natural de cada una (cuadratura y
  aproximación en \omterm{def:b3:topology:compact}{compactos}; análisis gaussiano; métodos minimax y de
  Fourier en cosenos). Una frase sobre qué dio la teoría general (Partes
  I y VI) que ningún cálculo individual podía dar.
\end{enumerate}

\medskip
\noindent\textbf{Parte VII --- El término de error y el núcleo que hay
detrás de los pesos.} Aquí $I$ es
\omterm{def:b3:topology:compact}{compacto} y $f \in \mathcal
C^{2n}(I)$.
\begin{enumerate}[resume]
  \item (Fórmula del error de Gauss) Sea $Hf$ el interpolante de Hermite
  de grado $\leq 2n - 1$ que iguala $f$ y $f'$ en los nodos
  $t_1, \dots, t_n$ (demuéstrense su existencia y el error puntual
        \[
        f(t) - Hf(t) =
        \frac{f^{(2n)}(\xi_t)}{(2n)!}\;p_n(t)^2
        \]
        mediante el habitual argumento de la función auxiliar).
        Dedúzcase, integrando esta identidad contra $w$ y encajando
        entre los extremos de $f^{(2n)}$, que
        \[
        \int_I f\,w - Q_n(f)
        = \frac{f^{(2n)}(\xi)}{(2n)!}\,h_n
        \qquad\text{para cierto } \xi \in I,
        \]
        con $h_n = \norm{p_n}^2$ como en la Parte VI: la cuadratura de
        Gauss se equivoca en una derivada $2n$-ésima, ponderada por la
        norma al cuadrado del polinomio ortogonal mónico. \item (Los
        pesos son valores de Christoffel) Usando el núcleo reproductor
        $K_n(x, y) = \sum_{k=0}^{n-1}
        \frac{p_k(x)p_k(y)}{h_k}$ de $\R_{n-1}[t]$ y la exactitud de
        $Q_n$ hasta el grado $2n - 2$, demostrar
        \[
        w_i \;=\;
        \Bigl(\,\sum_{k=0}^{n-1}
        \frac{p_k(t_i)^2}{h_k}\Bigr)^{\!-1} :
        \]
        cada peso es el valor en su nodo de la \emph{función de
        Christoffel} --- la positividad de los pesos (pregunta 10) otra
        vez, ahora con fórmula exacta. Compruébese que recupera
        $w_1 = w_2 = 1$ para $n =
        2$, $I = \intcc{-1}1$, $w = 1$. \item (Todo cuadra en una
        integral) Para el peso de Chebyshev y $n = 3$ nodos, calcúlense
        ambos miembros de
        \[
        \int_{-1}^{1}\frac{t^6}{\sqrt{1 - t^2}}\,\dd t
        = \frac{5\pi}{16},
        \qquad
        Q_3(t^6) = \frac{9\pi}{32},
        \]
        de modo que el error de la cuadratura es exactamente
        $\frac{\pi}{32}$; compruébese después que la fórmula del error
        de la pregunta 23 predice precisamente ese valor (aquí
        $f^{(6)} = 6!$ es constante y $h_3 = \norm{2^{-2}T_3}_w^2 =
        \frac\pi{32}$): la teoría y el cálculo concuerdan hasta la
        última cifra.
\end{enumerate}
\end{problem}
